\section{The common part and localized trace comparison}
\label{sec:traces}

The matrix model will now be placed in a Jiang--Su tensor factor by an
order-zero map.  Its expectation is still asymptotically the identity.
This produces a common part of the model and the original inclusion.
The remaining part is small in a sense that can be compared uniformly
over all lower semicontinuous extended traces, including traces that are
infinite away from an ideal.

\subsection{Order-zero maps into the Jiang--Su algebra}

Write $\tau_{\Z}$ for the unique tracial state of $\Z$.
We first construct order-zero matrix maps and positive contractions
with arbitrarily many orthogonal copies.

\begin{lemma}\label{supplies:order-zero}
For every integer $k\geq 1$ and every $\xi>0$ there is a c.p.c.
order-zero map $\iota\colon M_k\to\Z$ such that
\[
 \tau_{\Z}(\iota(1))>1-\xi,
 \qquad \|\iota(e_{11})\|=1.
\]
\end{lemma}

\begin{proof}
Choose a multiple $d\geq2$ of $k$ with $(d+1)^{-1}<\xi$, and fix a
unital embedding $\mu\colon M_k\to M_d$.  Put
$p=e_{d+1,d+1}\in M_{d+1}$.  In $M_d\otimes M_{d+1}$, the complement
of $1\otimes p$ is the tensor product of two $d$-dimensional spaces.
Let $V$ be the flip on this complement, extended by the identity on
the range of $1\otimes p$.  The unitary group of each finite matrix
block is path connected.  There is therefore a continuous unitary path
$U(t)$, $0\leq t\leq1$, from $1$ to $V$, all of whose members commute
with $1\otimes p$.

Define
\[
 \iota_0(x)(t)
   =U(t)\bigl(\mu(x)\otimes(1-tp)\bigr)U(t)^*.
\]
This is a c.p.c. order-zero map into $C([0,1],M_d\otimes M_{d+1})$.
Its endpoint values are
\[
 \iota_0(x)(0)=\mu(x)\otimes1,
 \qquad
 \iota_0(x)(1)=1\otimes\operatorname{diag}(\mu(x),0).
\]
Consequently its range lies in the prime dimension-drop algebra
\[
 I_{d,d+1}
 =\{f\in C([0,1],M_d\otimes M_{d+1}):
 f(0)\in M_d\otimes1,\ f(1)\in1\otimes M_{d+1}\}.
\]
The commutation of $U(t)$ with $1\otimes p$ gives the exact defect
\[
 1-\iota_0(1)=\bigl(t\mapsto t(1\otimes p)\bigr).
\]
For $1\leq r\leq d+1$, the functions
\[
 y_r(t)=\sqrt t\,(1\otimes e_{r,d+1})
\]
belong to $I_{d,d+1}$ and satisfy
\[
 y_r^*y_r=1-\iota_0(1),
 \qquad
 \sum_{r=1}^{d+1}y_ry_r^*=(t\mapsto t1)\leq1.
\]
Thus every tracial state $\sigma$ on $I_{d,d+1}$ satisfies
\[
 (d+1)\sigma(1-\iota_0(1))
   =\sigma\left(\sum_{r=1}^{d+1}y_ry_r^*\right)\leq1.
\]
The prime dimension-drop algebra $I_{d,d+1}$ embeds
unitally into $\Z$ by \cite[Theorem~2.1(i)]{Rordam}.
Composing such an embedding with
$\iota_0$ gives the trace estimate.  Finally,
$\|\iota_0(e_{11})(0)\|=1$, so isometry of the embedding gives the
asserted norm.
\end{proof}

\begin{lemma}\label{supplies:copies}
For every integer $m\geq2$ there are a positive contraction
$D_m\in\Z$ of norm one, contractions $v_{m,1},\ldots,v_{m,m}\in\Z$,
and a finite positive integer $K_m$ such that
\begin{equation}\label{supplies:copy-relations}
 v_{m,r}^*v_{m,s}=\delta_{rs}D_m,
 \qquad
 1=\sum_{l=1}^{K_m}y_{m,l}^*(D_m-1/2)_+y_{m,l}
\end{equation}
for suitable $y_{m,l}\in\Z$.
\end{lemma}

\begin{proof}
Take $\iota\colon M_m\to\Z$ from
Lemma~\ref{supplies:order-zero}, and set
\[
 D_m=\iota(e_{11}),\qquad v_{m,r}=\iota^{1/2}(e_{r1}).
\]
Order-zero functional calculus and the product identity
\eqref{prelim:order-zero} give the first identity in
\eqref{supplies:copy-relations}.  The element $(D_m-1/2)_+$ is
nonzero, hence full in the simple unital algebra $\Z$.
Lemma~\ref{prelim:full-compressions} supplies the second identity.
\end{proof}

We fix these choices once and for all.  No uniform bound on $K_m$ is
required.

\subsection{The model and its common part}

Throughout this section, $A$ is a nonzero separable unital nuclear
$C^*$-algebra.  Set
\begin{equation}\label{setup:ambient}
 R=A\otimes\Z_1\otimes\Z_2\otimes\Z_3,
 \qquad B=R_\omega,
 \qquad B_1=(A\otimes\Z_1)_\omega\subset B.
\end{equation}
The first Jiang--Su factor carries the model and its common part;
the second supplies positive contractions with arbitrarily many
orthogonal copies. The constant third factor is reserved for the
complementary weights in the unitary-equivalence argument.
The tensor-factor inclusions are unital.  Write
$\pi\colon A\to B_1\subset B$ for the constant first-factor
inclusion. Nuclearity and strong self-absorption of $\Z$
\cite[Examples~1.14]{TW} imply that
$R$ is nuclear, exact, and $\Z$-stable.

Use the tests $a_j,\gamma_j,x_j,b_j$ from \eqref{model:tests}.
The matrix-model construction, Lemma~\ref{model:construction},
gives finite-dimensional algebras $F_n$, powers
$k_n$ of two, and c.p.c. maps
\[
 A\xrightarrow{\alpha_n}F_n
   \xrightarrow{\beta_n}A\otimes M_{k_n},
\]
where $\beta_n$ is order zero, $\alpha_n$ is asymptotically order zero,
and $L_{k_n}\beta_n\alpha_n(a)\to a$ for every $a\in A$.
Denote the lemma's retained positive contractions by
$W_{n,j}\in A\otimes M_{k_n}$. For each fixed $j$ they satisfy
\begin{equation}\label{setup:model-retained}
 \|W_{n,j}-\beta_n\alpha_n(x_j)\|\longrightarrow0
\end{equation}
and
\begin{equation}\label{setup:model-traces}
 \tau(L_{k_n}(W_{n,j}))
   \leq\tau(x_j)+\delta_{n,j}d_\tau(b_j)
 \qquad\bigl(d_\tau(b_j)<\infty\bigr).
\end{equation}
Here $\tau$ ranges over all lower semicontinuous extended traces on
$A$, and the nonnegative constants $\delta_{n,j}$ are independent of
$\tau$ and tend to zero for fixed $j$.

Choose maps $\iota_n\colon M_{k_n}\to\Z_1$ from
Lemma~\ref{supplies:order-zero} with
$s_n:=\tau_{\Z}(\iota_n(1))\to1$.  Define
\begin{equation}\label{setup:transformed-model}
 \begin{aligned}
 T_n&=\id_A\otimes\iota_n,&
 \beta'_n&=T_n\beta_n,\\
 \theta_n&=\beta'_n\alpha_n,&
 w_{n,j}&=T_n(W_{n,j}).
 \end{aligned}
\end{equation}
These maps and elements take values in $A\otimes\Z_1$, which we
regard as a subalgebra of $R$.

The order-zero product identity for $\iota_n$ tensors with
$\id_A$, first on elementary tensors and then by continuity.  Thus
$T_n$ and $\beta'_n$ are c.p.c. order-zero maps.  Put
\[
 L=\id_A\otimes\tau_{\Z}\colon A\otimes\Z_1\longrightarrow A.
\]
The functional $\tau_{\Z}\circ\iota_n$ is a positive trace on
$M_{k_n}$.  Indeed, the square-root order-zero map satisfies
\[
 \iota_n^{1/2}(z)^*\iota_n^{1/2}(z)=\iota_n(z^*z),
 \qquad
 \iota_n^{1/2}(z)\iota_n^{1/2}(z)^*=\iota_n(zz^*).
\]
Its mass is $s_n$, so
\[
 LT_n=s_nL_{k_n}.
\]
It follows that $L\theta_n(a)\to a$ for every $a\in A$.
The product estimate for an order-zero map also gives
\[
 \|\theta_n(c)\theta_n(d)\|
 \leq\|\alpha_n(c)\alpha_n(d)\|\longrightarrow0
 \qquad(c,d\in A_+,\ cd=0).
\]
Consequently the maps $\theta_n$ induce a c.p.c. order-zero map
\begin{equation}\label{setup:theta}
 \theta\colon A\longrightarrow B_1,
 \qquad \theta(a)=(\theta_n(a))_\omega,
 \qquad (w_{n,j})_\omega=\theta(x_j).
\end{equation}
The same letter $L$ will denote the induced unital completely positive
map $B_1\to A_\omega$.  With the constant inclusion of $A$ into
$A_\omega$ understood, we have $L\theta=L\pi=\id_A$.

\begin{lemma}[The common part]\label{setup:common}
There is a positive contraction $e\in B_1$ such that, with $q=1-e$,
\begin{equation}\label{setup:common-equations}
 L(e)=0,\qquad
 [e,\pi(a)]=[e,\theta(a)]=0,\qquad
 q\theta(a)=q\pi(a)
 \quad(a\in A).
\end{equation}
In particular, if $h=\theta(1)$, then $qh=q$.
\end{lemma}

\begin{proof}
Both ranges $\pi(A)$ and $\theta(A)$ belong to the multiplicative
domain of $L$.  For $\theta$, the two Schwarz inequalities give
\[
 a^*a=L(\theta(a))^*L(\theta(a))
 \leq L(\theta(a)^*\theta(a))
 \leq L(\theta(a^*a))=a^*a.
\]
The analogous equalities hold with $a$ and $a^*$ interchanged.
The assertion for $\pi$ follows in the same way, or directly from
its multiplicativity.

Let $D_0=C^*(\pi(A),\theta(A))\subset B_1$.  This is a separable
unital algebra, $L|_{D_0}$ is a homomorphism, and
\[
 J=\ker(L|_{D_0})
\]
is an ideal containing $\theta(a)-\pi(a)$ for every $a\in A$.
Choose a positive contractive quasicentral approximate unit $(f_l)$
for $J$ in $D_0$.  Then $L(f_l)=0$, the commutators of $f_l$ with
both ranges tend to zero, and
\[
 (1-f_l)(\theta(a)-\pi(a))\longrightarrow0
 \qquad(a\in A).
\]
Lift the $f_l$ to positive contractive sequences in
$A\otimes\Z_1$.  Apply the countable diagonalization of
Lemma~\ref{prelim:epsilon-test} to the preceding equations on dense
countable subsets, together with the coordinate test
$\|L(f_{l,n})\|$.  The latter tends to zero along $\omega$ for each
fixed $l$.  The resulting positive contraction $e\in B_1$ satisfies
\eqref{setup:common-equations}.  Taking $a=1$ gives $qh=q$.
\end{proof}

Fix positive contractive representatives
$e_n\in A\otimes\Z_1$ of $e$, and write
\begin{equation}\label{setup:representatives}
 q_n=1-e_n,\qquad \ell_n=\|L(e_n)\|.
\end{equation}
Thus $\ell_n\to_\omega0$.  Exact vanishing of $L(e_n)$ at any
coordinate is neither asserted nor needed.

\subsection{Rank estimates on a prescribed ideal}

Fix $j$ temporarily, and abbreviate
\[
 a=a_j,\quad \gamma=\gamma_j,\quad
 x=x_j,\quad b=b_j,\quad w_n=w_{n,j}.
\]
Inside $A\otimes\Z_1\subset R$, set
\begin{equation}\label{trace:coordinate-positives}
 \begin{split}
 u_{n,0}&=x^{1/2}e_nx^{1/2},\\
 u_{n,1}&=w_n^{1/2}e_nw_n^{1/2},\\
 z_n&=a^{1/2}(q_n-1/2)_+a^{1/2}.
 \end{split}
\end{equation}
These are positive contractions.

\begin{lemma}[Localized rank estimates]\label{trace:rank}
For each $j$ there are nonnegative numbers
$\varepsilon_{n,j}\to_\omega0$, independent of the trace, with the
following properties.  Let $\lambda$ be any lower semicontinuous
extended trace on $R$, and let $\tau=\lambda|_A$.  If
$d_\tau(b)<\infty$, then, for $i=0,1$ and $\eta>0$,
\begin{equation}\label{trace:small-rank}
 d_\lambda((u_{n,i}-\eta)_+)
   \leq\eta^{-1}\varepsilon_{n,j}d_\tau(b).
\end{equation}
If $\ell_n\leq1/4$, then, with no finiteness assumption,
\begin{equation}\label{trace:large-rank}
 d_\lambda((z_n-\gamma/16)_+)
   \geq(\gamma/16)d_\tau(b).
\end{equation}
\end{lemma}

\begin{proof}
We first work under the hypothesis $d_\tau(b)<\infty$.
On the hereditary subalgebra
$\her_A(b)\otimes\Z_1$ one has
\begin{equation}\label{trace:localized-expectation}
 \lambda(y)=\tau(L(y))\qquad(y\geq0).
\end{equation}
To justify this formula, every positive contraction in
$\her_R(b\otimes1)$ has $\lambda$-value at most
$d_\lambda(b\otimes1)=d_\tau(b)$.  The restriction of $\lambda$
to $\her_A(b)\otimes\Z_1$ therefore extends to a bounded trace;
the same holds for $\tau$ on $\her_A(b)$.  For a positive
$s\in\her_A(b)$, the functional
\[
 v\longmapsto\lambda(s\otimes v)
\]
is a bounded positive trace on $\Z_1$, of mass $\tau(s)$.
Its trace identity follows by applying the trace identity for
$\lambda$ to $s^{1/2}\otimes z$.  Uniqueness of the tracial state
of $\Z$ gives $\lambda(s\otimes v)=\tau(s)\tau_{\Z}(v)$.
Linearity and density prove \eqref{trace:localized-expectation}.
In particular,
\begin{equation}\label{trace:first-small-part}
 \lambda(u_{n,0})
 =\tau(x^{1/2}L(e_n)x^{1/2})
 \leq\ell_n\tau(x)
 \leq\ell_n d_\tau(b).
\end{equation}

There is a separate estimate for $w_n$, which need not lie in that
hereditary subalgebra.  Define
\[
 \rho_n=\lambda\circ T_n,
 \qquad
 \sigma_n(s)=\rho_n(s\otimes1)\quad(s\in A_+).
\]
Order-zero square-root calculus shows that $\rho_n$ is an extended
trace; its lower semicontinuity follows from continuity of $T_n$.
Furthermore,
\[
 \sigma_n(s)
 =\lambda(s\otimes\iota_n(1))\leq\lambda(s\otimes1)=\tau(s).
\]
Lemma~\ref{prelim:matrix-traces}, applied to $\rho_n$, gives
\begin{equation}\label{trace:matrix-formula}
 \rho_n(Y)=\sigma_n(L_{k_n}(Y))
 \qquad(Y\in(A\otimes M_{k_n})_+).
\end{equation}
This identity includes infinite values. It follows from
\eqref{setup:model-traces} that
\begin{equation}\label{trace:retained-finite}
 \lambda(w_n)
 \leq\tau(L_{k_n}(W_{n,j}))
 \leq\tau(x)+\delta_{n,j}d_\tau(b)<\infty.
\end{equation}

Put
\[
 S_n=w_n^{1/2}q_nw_n^{1/2},\qquad
 S'_n=x^{1/2}q_nx^{1/2},\qquad
 \Delta_{n,j}=\|S_n-S'_n\|.
\]
The common-part identity implies
$\Delta_{n,j}\to_\omega0$: in the quotient, the two expressions
are respectively $q\theta(x)$ and $q\pi(x)$.
Compressing $S'_n\leq S_n+\Delta_{n,j}1$ by $x^{1/(2r)}$ gives
\[
 \lambda(x^{1/(2r)}S'_nx^{1/(2r)})
 \leq\lambda(S_n)+\Delta_{n,j}\tau(x^{1/r})
 \leq\lambda(S_n)+\Delta_{n,j}d_\tau(b).
\]
Here the compression of $S_n$ has trace at most $\lambda(S_n)$,
by traciality and contractivity of $x^{1/(2r)}$.
The compressed $S'_n$ converges in norm to $S'_n$ as $r\to\infty$.
Lower semicontinuity yields
\begin{equation}\label{trace:localized-norm-error}
 \lambda(S'_n)\leq\lambda(S_n)+\Delta_{n,j}d_\tau(b).
\end{equation}
All of $\lambda(w_n)$, $\lambda(S_n)$, $\lambda(S'_n)$, and
$\tau(x)$ are finite in the case under consideration.  We may
therefore subtract these trace values.  Since
$x=S'_n+u_{n,0}$ and $w_n=S_n+u_{n,1}$,
\eqref{trace:retained-finite} and
\eqref{trace:localized-norm-error} give
\[
 \lambda(u_{n,1})
 \leq\lambda(u_{n,0})
       +(\delta_{n,j}+\Delta_{n,j})d_\tau(b).
\]
Together with \eqref{trace:first-small-part}, this proves, for both
$i=0,1$,
\[
 \lambda(u_{n,i})
 \leq(\delta_{n,j}+\Delta_{n,j}+\ell_n)d_\tau(b).
\]
Set
\begin{equation}\label{trace:error-coefficient}
 \varepsilon_{n,j}=\delta_{n,j}+\Delta_{n,j}+\ell_n.
\end{equation}
The trace-cut estimate \eqref{prelim:trace-cut} proves
\eqref{trace:small-rank}.  In particular, when $d_\tau(b)=0$ the
right side is exactly zero, not an absolute error term.

We now prove \eqref{trace:large-rank}.  Suppose $\ell_n\leq1/4$ and
write
\[
 t_n=(z_n-\gamma/16)_+.
\]
Because $(q_n-1/2)_+\geq1/2-e_n$ and
$t_n\geq z_n-(\gamma/16)1$, we have
\begin{equation}\label{trace:expectation-lower}
 L(t_n)\geq(1/2-\ell_n)a-(\gamma/16)1.
\end{equation}
On the support of $b$, the element $a$ is at least $\gamma/2$.
Consequently compression by $b^{1/(2r)}$ gives
\begin{equation}\label{trace:compressed-lower}
 b^{1/(2r)}L(t_n)b^{1/(2r)}
 \geq(\gamma/16)b^{1/r}.
\end{equation}
If $d_\tau(b)<\infty$, the compression
$b^{1/(2r)}t_nb^{1/(2r)}$ belongs to
$\her_A(b)\otimes\Z_1$, so
\eqref{trace:localized-expectation} and traciality imply
\[
 \lambda(t_n)
 \geq\tau(b^{1/(2r)}L(t_n)b^{1/(2r)})
 \geq(\gamma/16)\tau(b^{1/r}).
\]
Let $r\to\infty$ and use
$d_\lambda(t_n)\geq\lambda(t_n)$, since $t_n$ is a contraction.
This proves \eqref{trace:large-rank} in the finite-rank case,
including rank zero.

For the infinite-rank case we use ideals, without appealing to
\eqref{trace:localized-expectation}.  First,
\begin{equation}\label{trace:expectation-ideal}
 L(t_n)\otimes1\in\Ideal_R(t_n).
\end{equation}
To see this, apply the domination theorem,
Theorem~\ref{prelim:domination}, on $\Z_1$ to
$v\mapsto\tau_{\Z}(v)1$ and the identity map.  The source is exact,
the maps are nuclear, the identity has order zero, and the required
pointwise ideal containment follows from simplicity of $\Z$.
There are therefore completely positive maps of the form
\[
 V_l(v)=\sum_s z_{l,s}^*vz_{l,s}
\]
converging pointwise in norm to $\tau_{\Z}(v)1$.  Including $1$ in
the approximation tests allows us to require $\|V_l\|\leq2$.
Tensoring with $\id_A$ then preserves pointwise norm convergence:
it holds first on elementary tensors, and the uniform norm bound
extends it to their completion.  Each resulting approximation to
$L(t_n)\otimes1$ is a sum of compressions of $t_n$ in $R$, proving
\eqref{trace:expectation-ideal}.

Let $I=\Ideal_R(t_n)$.  In $R/I$, equation
\eqref{trace:expectation-lower} and
\eqref{trace:expectation-ideal} give
\[
 0\geq(1/2-\ell_n)(a+I)-(\gamma/16)(1+I).
\]
Thus $\|a+I\|\leq\gamma/4$, and $(a-\gamma/4)_+\in I$.
Taking one more positive cut gives
\[
 b=\bigl((a-\gamma/4)_+-\gamma/4\bigr)_+,
 \qquad [b]\leq M[t_n]
\]
for some finite positive integer $M$.  If $d_\tau(b)=\infty$, then
$d_\lambda(b)=\infty$, so this comparison forces
$d_\lambda(t_n)=\infty$.  This finishes the proof in all cases.
\end{proof}

\subsection{Slowly chosen corners and bounded comparison witnesses}

We let the number of available orthogonal copies in the second
Jiang--Su factor grow slowly with the coordinate, so that the
localized comparisons remain valid.

\begin{proposition}[Localized comparison]\label{trace:comparison}
There is a sequence $m(n)\geq2$ with $m(n)\to_\omega\infty$ such
that, on putting
\[
 P_n=D_{m(n)}\in\Z_2,\qquad P=(P_n)_\omega\in B,
 \qquad
 \Psi(a)=P(q-1/2)_+\pi(a),
\]
the map $\Psi\colon A\to B$ is c.p.c. order zero and, for every
$j\geq1$,
\begin{equation}\label{trace:small-part-comparison}
 e\pi(x_j),\ e\theta(x_j)
 \precsim (\Psi(a_j)-c_j)_+
 \quad\text{in }B,
 \qquad c_j=\gamma_j/64.
\end{equation}
\end{proposition}

\begin{proof}
Temporarily fix $m\geq2$ and $j$, and keep the abbreviations from
the preceding proof.  The copy of $D_m$ is now in $\Z_2$, separate
from $z_n$ and $t_n=(z_n-\gamma/16)_+$.  The compression identity
in \eqref{supplies:copy-relations} gives
\begin{equation}\label{trace:separate-leg-comparison}
 [t_n]\leq K_m[t_n(D_m-1/2)_+]
       \leq K_m[(z_nD_m-\gamma/32)_+]
 \quad\text{in }W(R).
\end{equation}
For the second inequality, use joint functional calculus and the
scalar inequality
\[
 (s-\gamma/16)_+(t-1/2)_+
      \leq(st-\gamma/32)_+\qquad(0\leq s,t\leq1).
\]
Where the left side is positive, the difference between the
uncut expression on the right and the left side is
$s/2+(\gamma/16)t-\gamma/16\geq0$; elsewhere the assertion is
immediate.

Enumerate the countable set of pairs $(j,\eta)$ with
$j\geq1$ and $\eta\in\mathbb Q_{>0}$.  For each $m\geq2$ let
$\Omega_m$ be an $\omega$-large set on which $\ell_n\leq1/4$ and
\begin{equation}\label{trace:slow-requirement}
 \eta^{-1}\varepsilon_{n,j}
       \leq\frac{\gamma_j}{32K_m}
\end{equation}
for all the first $m$ pairs.  Such a set exists because the finitely
many error coefficients tend to zero along $\omega$.  Define
\[
 E_m=\{n:n\geq m\}\cap\bigcap_{r=2}^m\Omega_r.
\]
The sets $E_m$ are decreasing and $\omega$-large.  If $n$ belongs
to at least one $E_m$, let $m(n)$ be the largest such $m$; otherwise
put $m(n)=2$.  The largest value exists since membership in $E_m$
requires $m\leq n$.  This construction gives
$m(n)\to_\omega\infty$ and ensures that, for every fixed pair
$(j,\eta)$, condition \eqref{trace:slow-requirement} with
$m=m(n)$ holds on an $\omega$-large set.

Set $P_n=D_{m(n)}$.  For each fixed $(j,\eta)$, on the preceding
$\omega$-large set and for every lower semicontinuous extended
trace $\lambda$ on $R$, we claim
\begin{equation}\label{trace:uniform-gap}
 d_\lambda((u_{n,i}-\eta)_+)
 \leq\frac12 d_\lambda((z_nP_n-\gamma_j/32)_+)
 \qquad(i=0,1).
\end{equation}
When $d_\tau(b_j)<\infty$, combine
\eqref{trace:small-rank}, \eqref{trace:large-rank},
\eqref{trace:separate-leg-comparison}, and
\eqref{trace:slow-requirement}:
\[
 \begin{split}
 d_\lambda((u_{n,i}-\eta)_+)
 &\leq \frac{\gamma_j}{32K_{m(n)}}d_\tau(b_j)\\
 &\leq\frac12
      d_\lambda((z_nP_n-\gamma_j/32)_+).
 \end{split}
\]
This includes $d_\tau(b_j)=0$.  If $d_\tau(b_j)=\infty$, then
\eqref{trace:large-rank} gives $d_\lambda(t_n)=\infty$, and
\eqref{trace:separate-leg-comparison}, with finite $K_{m(n)}$,
forces the right side of \eqref{trace:uniform-gap} to be infinite.
Thus \eqref{trace:uniform-gap} also holds in that case.  Its
coordinate set is independent of $\lambda$.

The element $P$ commutes with the entire first subsystem $B_1$.
Also $q$ commutes with $\pi(A)$.  Hence $\Psi$ is a homomorphism
$\pi$ multiplied by commuting positive contractions, and is c.p.c.
order zero.  In $B$ we have
\begin{equation}\label{trace:coordinate-identifications}
 (u_{n,0})_\omega=e\pi(x_j),\qquad
 (u_{n,1})_\omega=e\theta(x_j),\qquad
 (z_nP_n)_\omega=\Psi(a_j).
\end{equation}

Since $R$ is exact and $\Z$-stable,
Lemma~\ref{prelim:half-comparison} applies to
\eqref{trace:uniform-gap}.  Thus, on an $\omega$-large set,
\[
 (u_{n,i}-\eta)_+
   \precsim(z_nP_n-2c_j)_+.
\]
We explain the passage to the ultrapower, because the norm bound
on witnesses is essential.  Fix $j$, a positive rational $\eta$,
$i\in\{0,1\}$, and an accuracy $\epsilon>0$, and let
$u_i=(u_{n,i})_\omega$.  By
Lemma~\ref{prelim:bounded-witness}, choose witnesses $r_n\in R$
on this large set such that
\[
 \begin{split}
 \bigl\|r_n^*(z_nP_n-c_j)_+r_n-(u_{n,i}-\eta)_+\bigr\|
     &<\epsilon,\\
 \|r_n\|^2&\leq(1+\epsilon)/c_j.
 \end{split}
\]
Set $r_n=0$ elsewhere.  This is a bounded sequence, so it defines
$r\in B$.  By \eqref{trace:coordinate-identifications},
\[
 \bigl\|r^*(\Psi(a_j)-c_j)_+r
       -(u_i-\eta)_+\bigr\|\leq\epsilon.
\]
As $\epsilon$ is arbitrary, each displayed source cut is Cuntz
below $(\Psi(a_j)-c_j)_+$.  Letting positive rational $\eta$
decrease to zero proves \eqref{trace:small-part-comparison}.
\end{proof}

\begin{lemma}\label{trace:copies}
For every positive integer $L$ there are contractions
$v_1,\ldots,v_L\in B$, represented entirely in the $\Z_2$ leg,
such that
\begin{equation}\label{trace:ultrapower-copies}
 v_l^*v_{l'}=\delta_{ll'}P\qquad(1\leq l,l'\leq L).
\end{equation}
In particular these elements commute with $B_1$, and hence with
$\pi(A)$ and $q$.
\end{lemma}

\begin{proof}
At coordinates where $m(n)\geq L$, use the elements
$v_{m(n),l}$ from Lemma~\ref{supplies:copies}; at all other
coordinates use zero.  The representatives have norm at most one.
Since $m(n)\to_\omega\infty$, their relations in
\eqref{supplies:copy-relations} give
\eqref{trace:ultrapower-copies}.  Their tensor-factor location gives
the stated commutation.
\end{proof}
